\section{Experimental Protocol and Reproducibility}
\label{sec:protocol}

\paragraph{Frozen rules.} Five rules governed every run, and we report violations
as well as compliance. (i) Thresholds are chosen on validation only. (ii) Every
evaluation that touches test is logged, numbered, and justified in advance.
(iii) Every number in this paper is produced by a committed script reading a
committed JSON or CSV. (iv) The original dataset and the pre-audit runs are never
modified. (v) Test labels are corrected only when objectively wrong, with
before/after numbers reported.

\paragraph{Test-touch accounting.} The full log contains $18$ numbered touches.
Four are large: the seed-0 frozen-protocol detector evaluation, the seed-0 and
seed-1/2 two-stage evaluations, and the corrected detector-aligned two-stage
re-scoring. The remainder are the per-seed detector evaluations, the MNv3
baseline evaluations, and the Grad-CAM passes, each sanctioned by the protocol
before execution. We also disclose a deviation: the upstream training script
auto-evaluates test at confidence $0.25$ at the end of training, and this
behaviour could not be disabled by flag. It is recorded as D017. No threshold,
checkpoint, or architecture choice references those numbers; they appear in no
table here. The frozen-protocol evaluation is a separate, single pass.

\paragraph{Correction we had to make.} An earlier version of the three-seed
two-stage result reused the seed-0 detector for all three seeds, producing a
tight mean of $0.879 \pm 0.018$. The arms A and B records were byte-identical
across seeds, which is the tell. Re-running the pipeline with each seed's own
detector gives $0.894 / 0.825 / 0.724$ and a standard deviation of $0.086$. The
tight interval was an artifact of a shared component. We report the corrected
numbers throughout and flag the earlier figure as withdrawn.

\paragraph{Compute.} All training and evaluation ran on one NVIDIA RTX~4070
Laptop GPU with $8$\,GB of memory, under Windows, within an approved budget of
roughly $30$ GPU-hours. Detector training took $50$--$85$ minutes per seed;
the second-stage classifier $30$ epochs per run; the ablation matrix five
classifier retrains. Evaluation uses $640\times640$ input, NMS IoU $0.7$,
\texttt{max\_det} $300$, FP32, no test-time augmentation. Latency figures are PC
GPU FP32 medians over $50$ repetitions after $10$ warm-up iterations, and are
\emph{not} edge-device numbers; no Jetson or TensorRT measurement was performed.

\paragraph{What the budget forbade.} We did not train YOLOv8m or YOLOv8l on the
clean split, did not evaluate a second detector family (RT-DETR or Faster
R-CNN), did not run grouped cross-validation, did not re-mine hard negatives on
data the detector had not seen, and did not retrain after the test-label
correction. Each omission is a real limitation and is listed in
Section~\ref{sec:limitations}. We prefer to name them than to let a single-seed
or single-family result stand in for a claim it cannot support.

\paragraph{Artifacts.} The repository contains the code, the six detector
checkpoints (three seeds, ONNX included), the three stage-2 checkpoints, the
frozen split manifest and its digest, all evaluation JSONs, the generated tables,
and the figures. Regenerating the tables requires no GPU and no dataset: the
committed JSONs are sufficient, and the table script fails loudly if a required
field is missing rather than silently substituting a value.
